\chapter{Conclusion}
\label{ch:conclusion}

ModFS shows how independently captured Ubuntu installations can be assembled over one exact shared base while retaining conventional package paths. The compatibility question requires several policies: package and module relations reject incompatible selections, allocation windows prevent numeric-identity collisions, and registry-specific reconciliation preserves state that ordinary layer precedence would discard. Runtime listeners and context-sensitive directory markers extend the conflict taxonomy beyond package metadata. The implemented registry inventory remains a defined scope, rather than proof that every maintainer-script effect has been covered.

The storage result is server-side and cohort-dependent. For the \dref{/catalogue/total}-module catalogue, the modelled monolithic-to-modular ratio is \dref{/storage/all/ratio}\(\times\), compared with \dref{/storage/small/ratio}\(\times\) for small modules and \dref{/storage/large/ratio}\(\times\) for large workloads. Sharing the base helps most when it would otherwise dominate each variant. The monolithic totals use an additive model calibrated on smaller modules; they do not establish equivalent savings for node transfers, deduplicating repositories, or every workload distribution.

Verification is the strongest methodological result. Admission accepts \dref{/tier1/totals/2/accept} pairs and \dref{/tier1/totals/3/accept} triples, while all \dref{/tier2/current/compose/n} sampled physical compositions pass the implemented predicates. Those observations coexist with directed counterexamples exposing omitted checks and real merger defects. Historical boots also show successful command probes alongside a failed service. A clean verdict establishes the predicate that produced it; broader correctness requires an observation capable of detecting the corresponding failure.

The physical sweep supports an approximately linear cost description within its sampled range. Similar composition slopes in successive sweeps are encouraging, but changes in inputs and environment prevent attributing the slope solely to the method. Likewise, archive pinning constrains package resolution without proving byte-identical builds. Digests bind existing artefacts; they neither remove generated state nor authenticate its publisher.

Future work should strengthen content-preservation checks, extend registry coverage, and collect matched composition and boot evidence from the same artefacts. Controlled independent rebuilds, large-workload monolithic baselines, and operation of the packaged GPU driver on matching hardware would address distinct remaining claims. Transactional publication, node updates, and rollback require a deployment mechanism and its own evaluation. The contribution established here is a bounded construction method, an explicit conflict taxonomy, and a clearer account of what its verification can and cannot observe.
